\documentclass{article}
\usepackage{amsmath,hyperref}
\title{Choosing tags across several possible source modes}
\date{}
\begin{document}
\maketitle
Suppose two scan stations can label each run by its physical epoch.
Both stations use the same fixed conversion from epoch to a shorter
tag before the observations are assembled. The tag is reliable, but
it may intentionally group epochs. A false join can swap records
within the reported tag, using a global allowance of $K$ displaced
identities. It cannot change, discard or duplicate readings.
We ask for the fewest tag values that let the controller choose the
right reserve location for every candidate mode and every allowed
join. There are only two possible actions. Modes recommending the
same action need not be distinguished.

\paragraph{Exact test for a proposed tag map.}
Use the finite-roster assumptions from
\href{https://github.com/eternal-november/eos-research-programmes/pull/405}{PR~405}:
both scan margins are fixed across modes within each epoch, with
known candidate $00$ counts $x^m_i$. A fixed tag partition $P$ induces
tagged counts $y_G^m=\sum_{i\in G}x_i^m$. Every tag becomes a block in
the join, even if its member epochs are not contiguous. For each
pair of opposite-action modes $m,m'$, check
\[
   \sum_{G\in P}|y_G^m-y_G^{m'}|>2\lfloor K/2\rfloor.
\]
This is necessary and sufficient: the observed tables must be
consistent only with modes prescribing one action. A common table
forces worst-case error at least one half, attainable by guessing.
The separate stronger requirement of preserving all known-mode
distances has a simpler conflict graph: two epochs may share a
tag only if they never have opposite strict signs in any pairwise
count difference. This follows from equality in the absolute-value
triangle inequality. The decision problem does not impose that
extra restriction. For example, a pair with differences $(3,-1,2)$
has atomic distance 6 and pooled distance 4; with $K=2$ pooling
still suffices.

\paragraph{Hardness for the decision problem itself.}
Given a nonempty simple graph on $n$ epochs, enumerate edges
$e=(u_e,v_e)$ by positive integers. Let $A=\lfloor K/2\rfloor+1$
and $s_e=(2A+1)e$. Create two modes for edge $e$. The first takes
action zero and has every count equal to $s_e$. The second takes
action one and adds $A$ to coordinate $u_e$ and subtracts $A$
from coordinate $v_e$. This is a valid finite-roster construction:
make each epoch have $2M$ records for any $M>\max_e s_e+A$.
A count $x$ then gives the clean table $(x,M-x,M-x,x)$, with
balanced margins and all cells positive. Rosters never need any
invalid counts or inconsistent underlying pairs.

Across distinct edges, the absolute pooled difference is
$n|s_e-s_f|>2\lfloor K/2\rfloor$. Aggregation by any partition
still separates those modes, since its distance is no less than
the absolute pooled difference. For the opposite-action pair
from the same edge, the tagged distance is zero if the endpoints
share a tag and $2A>2\lfloor K/2\rfloor$ otherwise. Hence the
minimum number of tags is exactly the input's chromatic number.
This proves NP-hardness for decision sufficiency even at fixed
$K=2$, not just distance preservation. A proposed tag map can be
verified by the polynomial count test above, yielding NP-completeness
for the bounded-tag decision problem.
\href{https://www.cs.umd.edu/~gasarch/BLOGPAPERS/Karp.pdf}{Karp}'s
graph-coloring hardness supplies the reduction premise. The reduction
assumes arbitrary reusable epoch labels, not a clock that can only
mark consecutive intervals. Many modes are necessary for this proof;
it does not show hardness with two candidate modes.

\paragraph{Consecutive clock buckets are different.}
Restrict the labels to consecutive intervals of epochs in the
graph-constructed family. A boundary after some index between the
endpoints of each edge is necessary and sufficient. Choose the
rightmost permissible boundary of the earliest-ending interval,
discard intervals it hits, and repeat. If an optimal solution uses
another boundary in the earliest-ending interval, replacing that
boundary by this rightmost one stays in all still-required intervals
it previously hit. This proves minimum bucket count. Three disjoint
edges between successive epoch pairs need four consecutive buckets
but only two reusable labels. This greedy rule is proved only for
the graph-constructed pair form, not arbitrary mode tables.

\paragraph{Two modes with any counts.}
For two opposite-action modes, two reusable tags always preserve
the full atomic distance: group the nonnegative count differences
under one tag and the negative differences under the other. Thus
one tag is optimal if the pooled distance clears the threshold; two
are optimal if pooling fails but the full distance clears it; no
tag map suffices if even the full distance fails.

For consecutive buckets, let $s_i$ be the prefix sum of the signed
differences and $F(k,i)$ the greatest distance obtainable by dividing
the first $i$ epochs into $k$ nonempty consecutive buckets. With
$F(0,0)=0$ and other impossible initial entries omitted,
\[
 F(k,i)=\max_{j<i}\{ F(k-1,j)+|s_i-s_j|\}.
\]
For each fixed $k$, a left-to-right sweep maintains the two prefix
maxima of $F(k-1,j)+s_j$ and $F(k-1,j)-s_j$. Subtract or add
$s_i$ respectively and take the larger. Thus the best distance
for every bucket count, and the minimum feasible count, take
$O(n^2)$ arithmetic operations. The first $k$ values take $O(kn)$.
A maximizing predecessor at each entry recovers a tag map. Exhaustive
cut enumeration checks the entire frontier on 1,024 signed inputs.

\paragraph{Scope.}
The result makes the tag allocation cost explicit when both stations
must implement it before the scan, without an oracle that observes the
mode. It does not price authentic tagging, certify clock alignment,
or establish that these many modes occur in a particular plant.
Ordinary graph coloring is the appropriate factored comparator on
the constructed family. The replay checks the characterizations
against exhaustive choices; it does not compare solver runtimes.
No historical priority, field calibration or full R3b closure is claimed.
\end{document}
